\scanpage{186}
\mspage{135}
lequel reçoit le diagramme non comm.
\diagram{\begin{tikzcd}[ampersand replacement=\&,row sep=small]
\&X_0\arrow[dl,"\alpha_1"']\arrow[dr,"\alpha_2"]\arrow[dd,dashed,"i_0", ""{name=C,near end}]\&\\
X_1\arrow[dr,"i_1"', ""{name=L,near end}]\&\&X_2\arrow[dl,"i_2", ""{name=R,near end}]\\
\&\widetilde{\mathscr X}\&{,}\\
\arrow[from=C,to=L,dashed,"\deleted{\unsure{f_0|i_0}}\widetilde\lambda_1"{swap}]\arrow[from=C,to=R,dashed,"\widetilde\lambda_2"]
\end{tikzcd}}
les flèches
\[
X_1\to\widetilde{\mathscr X}_{/1},\qquad X_2\to\widetilde{\mathscr X}_{/2}
\]
étant dans $W$. (et $=$ $W$-coasphériques -- elles ont $=$ un adjoint à gauche). Donc on trouve, comme précédemment
\[
\left\{\begin{aligned}
\alpha_1\in W&\Longleftrightarrow i_2\in W\\
\alpha_2\in W&\Longleftrightarrow i_1\in W ,
\end{aligned}\right.
\]
et de $=$ on peut \deleted{développer} formuler un \textquotedblleft lemme du cinq\textquotedblright\ comme le cor. page 125, mais avec $\widetilde{\mathscr X}$, \inserted{$\widetilde{\mathscr X}'$} au lieu de $\mathscr X,\mathscr X'$. \deleted{\unsure{De $=$, $\widetilde{\mathscr X}$}}

De plus $\widetilde{\mathscr X}$ est la solution d'un problème universel \inserted{dans $\mathit{Cat}$}, en termes de \unsure{données}
\diagram{\begin{tikzcd}[ampersand replacement=\&,row sep=small]
\&X_0\arrow[dl,"\alpha_1"']\arrow[dr,"\alpha_2"]\arrow[dd,"f_0", ""{name=C,near end}]\&\\
X_1\arrow[dr,"f_1"', ""{name=L,near end}]\&\&X_2\arrow[dl,"f_2", ""{name=R,near end}]\\
\&Y\&{,}\\
\arrow[from=C,to=L,"\widetilde\lambda_1"']\arrow[from=C,to=R,"\widetilde\lambda_2"]
\end{tikzcd}}
% archive sheet number: 186
